\usepackage{tcolorbox}
\usepackage{totcount,xfp}
\newtotcounter{totalpts}
%\setcounter{totalpts}{0} % Contador de puntos en tarea

\newcommand{\getTotal}{\number\totvalue{totalpts} }
\newcommand{\instruccion}{\begin{tcolorbox}[colback=blue!5!white,colframe=blue!75!black!70!white,title=Instrucciones]
Esta tarea es individual y se debe entregar en el buzón habilitado para eso en Canvas según los plazos definidos en la misma plataforma.  A menos que se indique lo contrario, todas las preguntas valen lo mismo. Para el desarrollo de los problemas, escriba el desarrollo en el mayor detalle posible, y defina todos los objetos matemáticos que use.

La tarea tiene \getTotal  puntos, pero el puntaje máximo será de \fpeval{floor({0.75*\getTotal})} puntos. En cada pregunta, el puntaje dado será: el total si está bien, la mitad si tiene un desarrollo concluyente pero tiene errores, y 0.0 puntos si tiene poco o nada de desarrollo.

Está permitido el uso de modelos de lenguaje (tipo chatGPT) para apoyar el desarrollo de la tarea, aunque sugerimos intentar un poco cada pregunta primero para generar dudas. Si bien no podemos controlar que le pidan directamente la respuesta al modelo, creemos es mejor para su proceso de aprendizaje que pidan sugerencias más que soluciones. Por esta razón, si deciden usarlo, debe adjuntar a su tarea otro pdf donde  muestren los prompts usados. Sugerimos usar prompts de este estilo: "Estoy desarrollando el siguiente ejercicio para una tarea en un curso introductorio de programación: [Copiar ejercicio]
  Y hemos visto hasta ahora los contenidos en las diapositivas adjuntas [adjunten los pdfs con las clases]. Dado este contenido, quería pedirte una sugerencia para resolver [explicar dificultad], ya que me confunde [escribir duda]. Como respuesta, no quiero que me des la solución, si no que una guía para ayudarme a iniciar el desarrollo de mi respuesta."

  Si tienen dudas y/o observaciones, por favor contáctenos por correo o Canvas. Es posible que el enunciado tenga errores. 
\end{tcolorbox}}


\newcommand{\code}[1]{\texttt{#1}}
% Logic: To create a counter with the points in each question and accumulate that in a total variable. Now I simply set the full threshold at the beginning to 75% and it gives the computed value automatically.
\newcommand{\pts}[1]{\if#11\textbf{[#1 punto]}\else\textbf{[#1 puntos]}\fi\addtocounter{totalpts}{#1}}
